\makeatletter
    \setlength{\absparsep}{18pt}
\makeatother

\begin{center}
    \textbf{\abstractitle}
\end{center}
\vspace{.5cm}

\noindent
Wave-induced dynamic loads on Tension Leg Platforms pose a threat to structural integrity, particularly due to the transmission of vibrations along the tendons. In this context, this study investigates the use of periodic configurations in these elements as a passive strategy for vibration mitigation. Based on Euler-Bernoulli beam theory, a numerical routine based on the Finite Element Method was developed and implemented in MATLAB, and validated against analytical formulation, to analyze the formation of band gaps, where wave propagation is suppressed. Two periodic configurations are evaluated: stepped beams (geometric periodicity) and bi-material beams (material periodicity). An analysis of band gaps formation is achieved via a periodic saturation parameter. The results demonstrate that, although both configurations are capable of generating band gaps, stepped beams exhibit higher tunability and yield broader attenuation bands. Thus, the optimization of geometric periodicity stands out as a robust design approach for filtering critical vibration frequencies in offshore systems.

\vspace{18pt}\noindent\textbf{Keywords}: Tension Leg Platforms, Periodic Beams, Vibration Mitigation, Bang Gaps, Finite Element Method.